\section{Related Work}
\label{sec:related_work}

\paragraph{Self-play with external verifiers.}
Self-play has long improved agents in closed games~\citep{silver2017mastering,silver2018general}, inspiring recent language-model work on alignment, instruction following, and reasoning through self-generated data or verifiable rewards~\citep{cheng2024self,chen2024self,cheng2024spar,huang2025spotlight,kuba2025language,ren2026seif}. Recent self-evolving systems often obtain reliable feedback by grounding task generation in external environments. For example, Absolute Zero~\citep{zhao2026absolute} and verified-code pipelines~\citep{wilf2025propose} rely on code executors; SPIRAL and MARSHAL~\citep{liu2025spiral,yuan2025marshal} use game simulators; and Search Self-play~\citep{lu2025search} leverages retrieval augmentation. While these works demonstrate the value of reliable verification, their specialized environments inherently bottleneck the task distribution, confining the agent's learning to narrow, predefined operational domains.

\paragraph{Synthetic task generation and filtering.}
To broaden the task space, general self-improvement pipelines commonly combine synthetic task generation with automated verification~\citep{yang2026self}. Representative mechanisms include unit tests for coding tasks~\citep{jimenez2024swe,gehring2024rlef,jiang2025coderl+}, success signals in interactive web environments~\citep{liu2024agentbench,zhou2024webarena,yao2024tau,song2026envscaler}, schema checks for tool use~\citep{li2023api,patil2025berkeley}, and deterministic constraint checkers for reasoning~\citep{zhou2023instruction,jiang2024followbench,qin2024infobench,wen2024benchmarking,pyatkin2025generalizing}. Although these mechanisms make synthetic data auditable, they primarily operate as passive, \emph{post-hoc} filters. They dictate which tasks survive but provide no proactive structural guidance to the proposer, causing unguided generation to frequently collapse into ill-posed tasks or redundant templates.

\paragraph{Skill interfaces for agents.}
Recent work studies agent skills as reusable procedural interfaces, with an emphasis on retrieval, compression, and progressive disclosure~\citep{ling2026agent,jiang2026sok,cho2026skillret,li2026skillsinjector,huang2026persistent,gao2026skillreducer,chen2026skilljuror,chen2026skill,he2026gems}. Another line of research automatically constructs or evolves skills from interaction traces and execution feedback~\citep{mi2026procmem,liu2026harnessing,wang2026skillx,shen2026skillfoundry,yang2026skillmaster,liu2026skillrevise,gautam2026skillaxe,wang2026skillgrad,liu2026skillsvote,ni2026trace2skill}, although preserving the fidelity of automatically generated skills remains challenging~\citep{li2026skillsbench,zhong2026skilllearnbench,zhou2026skillgenbench}. Most prior work uses skills to support inference-time execution, while a few recent studies explore their role during training~\citep{lu2026skill0}. In contrast, Skill-SP uses an evolving skill library as a training-time interface for task synthesis, enabling structural guidance and verification to co-evolve with the self-play curriculum. The following section formalizes this framework.